\documentclass[10pt]{article}
\usepackage[UTF8]{ctex}
\usepackage{amsmath, amssymb}
\usepackage{geometry}
\geometry{a4paper, margin=1.5cm}
\usepackage{enumitem}

\title{导数应用：从变化率到函数分析——写给高二文科生的进阶笔记}
\author{}
\date{}

\begin{document}
\maketitle

\section*{致亲爱的同学}
如果你已经掌握了导数的定义和基本求导法则，那么恭喜你，你已经拥有了一把强大的工具。接下来，我们要学习的是：\textbf{用这把工具去解决实际问题}。

在高考中，导数几乎从不单独以“求导”的形式出现，而是作为手段，去求切线、判断单调性、求极值与最值、解决优化问题等。本笔记将带你逐一攻克这些应用场景。

\section{切线问题——导数最直观的几何意义}
\subsection{切线的本质回顾}
我们知道，函数 $y=f(x)$ 在点 $(x_0, f(x_0))$ 处的导数 $f'(x_0)$，就是曲线在该点处\textbf{切线的斜率}。因此，切线方程为：
$$y - f(x_0) = f'(x_0)(x - x_0).$$
\begin{quote}
    \textbf{关键：} 切点必须在曲线上。
\end{quote}

\subsection{已知切点求切线}
\textbf{步骤：}
\begin{enumerate}
    \item 求导 $f'(x)$；
    \item 代入 $x_0$ 得斜率 $k = f'(x_0)$；
    \item 用点斜式写出切线方程。
\end{enumerate}

\textbf{例 1：} 求曲线 $y = x^3 - 2x + 1$ 在点 $(1,0)$ 处的切线方程。

\textbf{解：} $f'(x) = 3x^2 - 2$，$f'(1) = 1$，切线方程 $y - 0 = 1 \cdot (x - 1)$，即 $y = x - 1$。

\textbf{易错点：} 必须先验证点是否在曲线上。若题目说“在点 $(1,0)$ 处”，通常默认点在曲线上，但最好代入检查：$1^3 - 2\times1 + 1 = 0$，正确。

\subsection{过曲线外一点求切线（设切点法）}
如果点 $(a,b)$ 不在曲线上，就不能直接代入。此时需要\textbf{设切点}，列方程求解。

\textbf{标准解法：}
\begin{enumerate}
    \item 设切点为 $(x_0, f(x_0))$；
    \item 写出切线方程：$y - f(x_0) = f'(x_0)(x - x_0)$；
    \item 将已知点 $(a,b)$ 代入，得到关于 $x_0$ 的方程；
    \item 解出 $x_0$（可能有多个解，对应多条切线）；
    \item 代回得切线方程。
\end{enumerate}

\textbf{例 2：} 求过点 $(0, -4)$ 且与曲线 $y = x^2$ 相切的直线方程。

\textbf{解：} 设切点 $(x_0, x_0^2)$，则斜率 $k = 2x_0$。

切线方程：$y - x_0^2 = 2x_0 (x - x_0)$。

代入 $(0, -4)$：$-4 - x_0^2 = 2x_0 (0 - x_0) = -2x_0^2$。

整理：$-4 - x_0^2 = -2x_0^2 \Rightarrow x_0^2 = 4 \Rightarrow x_0 = \pm 2$。

当 $x_0 = 2$ 时，切点 $(2,4)$，$k=4$，切线 $y-4=4(x-2) \Rightarrow y=4x-4$；

当 $x_0 = -2$ 时，切点 $(-2,4)$，$k=-4$，切线 $y-4=-4(x+2) \Rightarrow y=-4x-4$。

\textbf{答案：} $y=4x-4$ 或 $y=-4x-4$。

\textbf{易错点：} 设切点时，$f(x_0)$ 要用原函数表达式，不能直接用 $y$；解方程后记得验证 $x_0$ 是否合理（如分母不为零等）。

\section{单调性——导数符号的“天气预报”}
\subsection{导数与单调性的关系}
若在某个区间内：
\begin{itemize}
    \item $f'(x) > 0$，则 $f(x)$ \textbf{单调递增}；
    \item $f'(x) < 0$，则 $f(x)$ \textbf{单调递减}；
    \item $f'(x) = 0$ 的点可能是极值点，也可能是拐点或平台。
\end{itemize}
\textbf{直观理解：} 导数表示瞬时变化率。正数说明函数值在增加，负数说明在减少。

\subsection{求单调区间的标准步骤}
\begin{enumerate}
    \item \textbf{确定定义域}（尤其对数、分式、根式函数）；
    \item \textbf{求导} $f'(x)$；
    \item \textbf{解方程} $f'(x) = 0$，得到临界点；
    \item \textbf{列表}：将定义域用临界点分成若干区间，判断每个区间内 $f'(x)$ 的符号；
    \item \textbf{写出单调区间}。
\end{enumerate}

\textbf{例 3：} 求函数 $f(x) = x^3 - 3x^2 + 1$ 的单调区间。

\textbf{解：} 定义域 $\mathbb{R}$。$f'(x) = 3x^2 - 6x = 3x(x-2)$。

令 $f'(x)=0$，得 $x=0$ 或 $x=2$。

列表：
\begin{itemize}
    \item 区间 $(-\infty,0)$：$f'(x)>0$，$f(x)$ 递增；
    \item 点 $0$：$f'(x)=0$，对应极大值；
    \item 区间 $(0,2)$：$f'(x)<0$，$f(x)$ 递减；
    \item 点 $2$：$f'(x)=0$，对应极小值；
    \item 区间 $(2,+\infty)$：$f'(x)>0$，$f(x)$ 递增。
\end{itemize}
递增区间：$(-\infty,0)$ 和 $(2,+\infty)$；递减区间：$(0,2)$。

\textbf{易错点：}
\begin{itemize}
    \item 不要忘记定义域；
    \item 单调区间之间用“和”连接，不能用“$\cup$”；
    \item 导数等于0的点要单独列出，但单调区间通常写成开区间。
\end{itemize}

\section{极值与最值——寻找函数的“山峰”与“谷底”}
\subsection{极值的定义与判定}
\begin{itemize}
    \item \textbf{极大值}：函数在某点附近，该点函数值最大（局部）；\textbf{极小值}：附近最小。
    \item \textbf{必要条件}：极值点处 $f'(x)=0$ 或导数不存在。
    \item \textbf{充分条件（一阶导数法）}：设 $x_0$ 是 $f'(x)=0$ 的点，若在 $x_0$ 左侧 $f'(x)>0$，右侧 $f'(x)<0$，则 $x_0$ 是极大值点；若左负右正，则是极小值点。
\end{itemize}

\subsection{求极值的步骤}
\begin{enumerate}
    \item 求定义域，求导；
    \item 解 $f'(x)=0$，得到所有临界点；
    \item 列表，判断每个临界点两侧导数符号；
    \item 写出极值点及极值。
\end{enumerate}

\subsection{闭区间上的最值}
最值可能出现在\textbf{端点}或\textbf{临界点}。

\textbf{步骤：}
\begin{enumerate}
    \item 求导，找出区间内的临界点；
    \item 计算端点函数值和各临界点函数值；
    \item 比较，最大值和最小值即为所求。
\end{enumerate}

\textbf{例 4：} 求 $f(x)=x^3-6x^2+9x+2$ 在 $[-1,4]$ 上的最大值与最小值。

\textbf{解：} $f'(x)=3x^2-12x+9=3(x-1)(x-3)$。

临界点 $x=1$，$x=3$（均在区间内）。

$f(-1)=-14$，$f(1)=6$，$f(3)=2$，$f(4)=6$。

最大值 $6$（在 $x=1$ 和 $x=4$ 处），最小值 $-14$（在 $x=-1$ 处）。

\textbf{易错点：} 最值可能出现在端点，务必计算；若区间是开区间，最值可能取不到。

\section{实际优化问题——导数的现实力量}
优化问题是高考实际应用题的热点，通常需要先建立函数模型，再求最值。

\subsection{建立函数模型}
从几何问题（如容积最大、面积最大）或实际情境（成本最低、利润最大）中抽象出目标函数，并明确自变量的取值范围（定义域）。

\subsection{解题流程}
\begin{enumerate}
    \item \textbf{设变量}：设未知量（如边长、时间等），通常一个变量即可；
    \item \textbf{列函数}：根据几何关系或实际关系写出目标函数 $f(x)$；
    \item \textbf{确定定义域}：根据实际意义确定 $x$ 的范围；
    \item \textbf{求导，找临界点}；
    \item \textbf{比较端点与临界点}，得出最值；
    \item \textbf{回答实际问题}（带单位）。
\end{enumerate}

\textbf{例 5：} 用一块长为 $8$ 米，宽为 $5$ 米的矩形铁皮，在四个角剪去四个相同的小正方形，然后折成一个无盖长方体容器。问剪去的小正方形边长为多少时，容器的容积最大？最大容积是多少？

\textbf{解：} 设剪去的小正方形边长为 $x$ 米，则 $0 < x < 2.5$。

折后长方体：长 $8-2x$，宽 $5-2x$，高 $x$。

容积 $V(x) = (8-2x)(5-2x)x = 4x^3 - 26x^2 + 40x$。

求导：$V'(x)=12x^2-52x+40=4(3x^2-13x+10)$。

令 $V'(x)=0$，解得 $x=1$ 或 $x=\frac{10}{3}\approx3.33$（舍去，超出定义域）。

计算端点：$V(0)=0$，$V(2.5)=0$，$V(1)=18$。

所以当 $x=1$ 米时，容积最大为 $18$ 立方米。

\textbf{易错点：} 定义域很容易被忽略，尤其是当 $x$ 过大导致长或宽为负时，必须限制。

\section{含参讨论——分类思想的训练}
参数出现在函数中时，导数的符号可能随参数变化，需要分类讨论。

\subsection{含参单调性讨论}
典型问题：已知函数在 $\mathbb{R}$ 上单调递增，求参数范围。

\textbf{方法：} $f'(x) \ge 0$ 恒成立（且不恒为0），转化为二次函数恒成立问题。

\textbf{例 6：} 已知 $f(x)=\frac{1}{3}x^3 - a x^2 + 3x + 1$ 在 $\mathbb{R}$ 上单调递增，求 $a$ 的取值范围。

\textbf{解：} $f'(x)=x^2 - 2a x + 3$。

要求 $f'(x) \ge 0$ 对一切 $x\in\mathbb{R}$ 恒成立。

二次函数开口向上，只需判别式 $\Delta \le 0$：
$$\Delta = 4a^2 - 12 \le 0 \Rightarrow a^2 \le 3 \Rightarrow -\sqrt{3} \le a \le \sqrt{3}.$$
\textbf{答案：} $a \in [-\sqrt{3}, \sqrt{3}]$。

若函数在某个区间上单调递增，则需在该区间内 $f'(x) \ge 0$，可能转化为最值问题。

\subsection{含参极值讨论}
当 $f'(x)=0$ 的解含参数时，需根据参数变化讨论极值点的个数和性质。

\textbf{例 7：} 讨论 $f(x)=x^3 - 3ax + 1$ 的极值情况。

\textbf{解：} $f'(x)=3x^2 - 3a = 3(x^2 - a)$。
\begin{itemize}
    \item 若 $a \le 0$，则 $f'(x) \ge 0$ 恒成立，$f(x)$ 在 $\mathbb{R}$ 上单调递增，无极值；
    \item 若 $a > 0$，则 $f'(x)=0$ 得 $x = \pm\sqrt{a}$。
    列表：
    \begin{itemize}
        \item 区间 $(-\infty,-\sqrt{a})$：$f'(x)>0$，$f(x)$ 递增；
        \item 点 $-\sqrt{a}$：$f'(x)=0$，对应极大值；
        \item 区间 $(-\sqrt{a},\sqrt{a})$：$f'(x)<0$，$f(x)$ 递减；
        \item 点 $\sqrt{a}$：$f'(x)=0$，对应极小值；
        \item 区间 $(\sqrt{a},+\infty)$：$f'(x)>0$，$f(x)$ 递增。
    \end{itemize}
    极大值 $f(-\sqrt{a})=2a\sqrt{a}+1$，极小值 $f(\sqrt{a})=-2a\sqrt{a}+1$。
\end{itemize}

\textbf{易错点：} 讨论时要注意参数是否影响定义域，以及二次项系数是否为零等特殊情形。

\section{综合应用（选讲）}
\subsection{由导函数图像推断原函数性质}
若给出 $y=f'(x)$ 的图像，可以判断 $f(x)$ 的单调区间和极值点：
\begin{itemize}
    \item $f'(x) > 0$ 的部分对应 $f(x)$ 递增；
    \item $f'(x) = 0$ 且两侧符号改变的点对应极值点。
\end{itemize}

\subsection{利用导数证明不等式}
通常将不等式转化为函数最值问题。例如证明 $e^x \ge x+1$，可设 $g(x)=e^x - x -1$，求导得 $g'(x)=e^x-1$，分析单调性，得 $g(x) \ge g(0)=0$。

\subsection{函数零点与导数}
结合单调性和极值的符号，可以判断函数零点的个数。例如三次函数 $f(x)=x^3-3x+1$，通过求极值并观察极值正负，可确定零点个数。

\section*{附录：高考常见题型速查表}
\begin{tabular}{|p{5cm}|p{8cm}|}
\hline
\textbf{题型} & \textbf{易错点提醒} \\
\hline
切线方程（点在曲线上） & 直接求导代值；易错点是忘记验证点是否在曲线上。 \\
\hline
切线方程（点不在曲线上） & 设切点列方程；易错点是忽略有多条切线的情况。 \\
\hline
求单调区间 & 解 $f'(x) > 0$ 和 $<0$；易错点是忘记定义域。 \\
\hline
求极值 & 列表判断导数符号；易错点是混淆极值与最值。 \\
\hline
闭区间最值 & 比较端点与临界点；易错点是漏算端点。 \\
\hline
含参单调性 & 二次函数恒成立问题；易错点是忘记讨论二次项系数为0。 \\
\hline
实际优化 & 建立函数，确定定义域；易错点是定义域错误导致最值无效。 \\
\hline
不等式证明 & 构造函数，求最值；易错点是忽略等号成立条件。 \\
\hline
\end{tabular}

\section*{结语}
希望这份笔记能帮助你从“会求导”走向“会用导数”。导数是一座桥梁，连接了代数与几何、数学与现实。多动手练习，你一定能掌握它！

\textbf{建议：} 先独立完成每个例题，再对照解答。遇到含参问题，多问自己“参数变化时，什么会改变？”祝你学习顺利！

\end{document}